\documentclass[12pt]{book}
\usepackage[utf8]{inputenc}
\usepackage[T1]{fontenc}
\usepackage{amsmath,amssymb,amsfonts}
\usepackage{amsthm}

\begin{document}

\setcounter{page}{342}
1) the \(b_{i}\) generate
B as an A-algebra, and
2) space. the elements \(b_{i} \otimes 1\) generate B~K as a K-vector

Then for each i,J, \(b_{i} b_{j} \otimes 1 \in B \otimes K\) , so we can write
\[b_{i} b_{j} \otimes 1=\sum_{k=1}^{n} b_{k} \otimes \lambda_{i j k}, \lambda_{i j k} \in k .\]

Let f be a common denominator for all \(\lambda_{i j k}\) , \(1 ≤i, j\) , \(k ≤n\) .Then the \(\lambda_{i j k}\) are all in \(A_{f}\) , so the \(A_{f}\) -module \(B_{0}\) generated by the \(b_{i} \otimes 1\) is in fact a ring. But \(B_{0}\) contains a set of generators for \(B \otimes A_{f}\) as an \(A_{f}\) -algebra, so \(B_{0}=B \otimes A_{f}\) , which shows that the latter is a finitely generated \(A_{f}\) -module.
q.e.d.

\setcounter{page}{343}
Lemma 4.4. Let \(f: X \to Y\) be a morphism, and let \(K^{*}\) be a residual complex on Y. Then there is a unique map
\[Tr_{f}: f_{*} f^{\Delta} K^{*} \to K^{*}\]
of graded sheaves on Y , such that whenever u is an open subset of Y, and W is a closed subscheme of \(f^{-1}(u)\) ,finite over U ,we have a commutative diagram

(Here we denote by f the restriction of f to \(f^{-1}(u)\) ,by \(K^{*}\) the restriction of \(K^{*}\) to v , and so forth.)

\setcounter{page}{344}
Proof. Let d denote the codimension function on Y associated with the residual complex \(K^{*}\) , and let it also denote the codimension function on X associated with the residual complex \(f^{\Delta} K^{*}\) (the one we mean will always be clear from the context).

Then according to the definition of the residual complexes we have isomorphisms
\[K^{p} \cong \sum_{d(y)=p} J(y)\]
and
\[\left(f^{\Delta} K^{\cdot}\right)^{p} \cong \sum_{d(x)=p} J(x) .\]

We will denote by \(J(y)\) (resp. \(J(x))\) the unique subsheaf of \(K^{p}\)
(resp. \(f^{\Delta}(K^{\cdot})^{p})\) which is isomorphic to an injective hull of \(k(y)\) (resp. \(k(x))\) , so as to make these isomorphisms canonical.

\setcounter{page}{345}
I claim that to give a map of graded \(\sheaf{Y}\) -modules
\[Tr_{f}: f_{*} f^{\Delta} K^{*} \to K^{*}\]
is equivalent to giving, for each \(x \in X\) which is closed in its fibre, a map
\[Tr_{f, x}: f_{*} J(x) \to J(y)\]
where \(y=f(x)\) .Indeed, x is closed in its fibre if and only if tr.d. \(k(x) / k(y)=0\) , i.e., if and only if \(d(x)=d(y)\) ,by
Proposition 3.4 above. Hence \(Tr_{f}\) certainly gives rise to a collection of maps \(Tr_{f, x}\) as above.

On the other hand, if \(d(y') \ge d(y)\) ,then there is no non-zero map of \(f_{*} J(x)\) into \(J(y')\) ,because \(y' \notin \overline{\{y\}}\) .Thus \(Tr_{f}\) is determined by the maps \(Tr_{f, x}\) above, and these can be given arbitrarily.

We will now construct maps \(Tr_{f, x}\) as above for each point \(x \in X\) which is closed in its fibre.

\setcounter{page}{346}
Given an \(x \in X\) closed in its fibre, choose by Lemma 4.3 an open neighborhood v of \(y=f(x)\) such that \(Z \cap f^{-1}(v)\) is finite over V. Let I be the ideal of Z of Z with the reduced induced structure, and let \(z_{n}\) , \(n=1,2, \cdots\) , be the subscheme of
\(I^{n}\) .For \(n<n'\) we have a closed immersion \(j_{n n'}: z_{n} \to z_{n'}\) , and one can verify (:) using VAR l, VAR 2, and (viii) above, that the following diagram is commutative:

One can also write down another diagram with three indices n , \(n'\) and \(n "\) which shows that the sheaves \(f_{*} i_{n *} i_{n}^{Y} f_{n}\) and the maps \(\rho_{j_{n n'}}\) form a direct system as n varies, and this diagram shows that the \(\rho_{i_{n}}\) map this direct system into \(f_{*} f^{\Delta}(K^{*})\) and the \(\rho_{g_{n}} \psi_{g_{n}}\) nap it to \(K^{*}\) . Furthermore the \(\rho_{i_{n}}\) are all injective maps. Thus we can pass to the limit and obtain a map of a certain subsheaf of \(f_{*} f^{\Delta}(K^{*})\) to \(K^{*}\) . Looking at the effect of this construction on the component \(f_{*} J(x)\) of \(f_{*} f^{\Delta} K^{*}\) , we see that the \(n^{th }\) term of the direct system, via the inclusion \(\rho_{i_{n}}\) ,is just \(f_{*} \underline{Hom}_{\sheaf{X}}(\sheaf{X} / I^{n}, J(x))\) , and hence the limit is \(f_{*} J(x)\) itself, since \(J(x)\) has support on Z Hence the map we obtain is defined on all of \(f_{*} J(x)\) , and we decree this to be \(Tr_{f, x}\)

\setcounter{page}{347}
To complete the proof of the lemma, we must check various things. First, that \(Tr_{f, x}\) is well-defined, i.e., does not depend on the choice of the open set V. That is clear, since all of our constructions are compatible with localization.

Second, we must check the property of the lemma. So let w and U be as in the statement. It will be sufficient to check the diagram on the component of x for each \(x \in W\) . Let Z be the closure of x , with the reduced induced structure (which will be finite over U). Let J be the ideal of z in W , and for each \(n=1,2,\) ,let \(z_{n}\) be as above. Then there are closed immersions as shown. We leave to the reader (') to write down a huge commutative diagram of isomorphisms which in the limit gives the diagram of the lemma on the component of x.

\setcounter{page}{348}
Proof of theorem. For a morphism f: \(X \to Y\) , we define \(Tr_{f}\) to be the map given by the lemma. It is clearly functorial in \(K^{*}\) In case f is a finite morphism, we can take U ~ Y and \(W=X\) , so that \(Tr_{f}=\rho_{f} \psi_{f}\) , which proves condition TRA 2.

To prove TRA i, let \(f: X \to Y\) and \(g: Y \to Z\) be two morphisms. It is sufficient to prove the condition of TRA 1 for a single residual complex \(K^{*}\) on Z and for a single \(x \in X\) which is closed in its fibre over Z. The question is local, so we may assume that \(w=\overline{\{x\}}\) is finite over z , and that \(v=\overline{\{y\}}\) , where \(y=f(x)\) , is finite over Z Finally it is sufficient to check the commutativity for a given element \(a \in \Gamma(J(x))\) . We choose a subscheme structure \(w_{n}\) on W with n large enough so that a is in the image of \(_{i}^{Y} J(x)\) .Now one can check (~) the required commutativity using the property of the lemma, and (viii), VAR 1,
and VAR 2 above.

The uniqueness of \(Tr_{f}\) is clear, as it was in the lemma.
q.e.d.

\end{document}